\typesetchapter{Chapter 1}{Ascension}


% 1-The Hall of Ceremony is introduced as a mined vertical cone: an impossible upward hollow cut into the Mountain to house the pendulum

The Hall of Ceremony had been prepared for three days, that was the official sentence, prepared, the word they used for that whole patient 
business of lowering lamps, filling bowls, polishing dull stone into a greasy obedience, brushing grit from the floor, tightening the posts 
of the low barrier, checking the brass pins, the cords, the little hinges, the miserable practical bones of reverence; and all this not for 
beauty, not first, but because the Pendulum was beautiful only from the proper distance, because under it beauty became weight, edge, 
passage, return, and a boy’s skull would have opened there as easily as a dropped jar.\par

No, the Hall had not been prepared in three days. It had merely been cleaned for us. The true preparation had taken centuries 
and left its illness in the stone: a huge upward hollow mined into the Mountain, round at the floor, narrowing as it climbed, 
not quite a cone, never quite anything one could name without lying a little, with walls that seemed to have been scraped, corrected, 
worried inward by generations of hands, and high above, far past the pillars and their dark ring of support, the upper throat from 
which the Thread descended, thin and impossible, holding at its end the disk that crossed the empty circle with the calm of a thing 
too heavy to hate us.\par

% 2- Five Pillars, I'm from light.

Below that falling Thread, the pillars had been set in a ring around the emptiness, five enormous things, too thick, too worked, too crowded with carvings, 
as if the Mountain had suffered a disease of symbols and the stone had swollen with them. They did not rise nobly, whatever the teachers said. 
They squatted first, they shouldered upward, they dragged their engraved bodies toward the dark annular band where they stopped, 
and every one of them seemed exhausted by having to represent more than stone should ever be asked to represent. 
Their bases were rubbed smooth by processions. 
Their lower cuts were blackened by hands, oil, dust, the old grease of reverence. 
Higher up, where no hand could reach, the carvings sharpened and multiplied, smug in the lamplight, clean because only the eye was allowed to touch them. 
Above, ribs of stone ran from pillar to pillar and made the star we had scratched since childhood in dirt and condensation, only here it had lost all 
playfulness, all quickness, and become a hard, fixed, oppressive geometry, a child’s sign fattened into architecture and hung over the place where no child was allowed to step.\par

On the lower side of the Hall, between Shield and Balance, the Mountain gates sat shut under the Genesis gravure, 
squat doors under a crowded scripture, the whole arrangement ugly with usefulness. 
Shield’s pillar stood to one side of that axis, full of bars, gates, narrow slits, hands pressed against closed stone, 
faces turned away from the ceremony as if the true event were always behind us, in the passages, where something might enter. 
Balance stood to the other side, more crowded still, with carved channels, jars, roots, grain steps, shells, little curls rising from stone water, 
all of it dry and yet somehow damp to the eye, like a pantry dreaming of a spring. 
Then Faith, with its waves or songs or whatever they were meant to be, lines climbing, bending, returning, mouths without faces, hands without wrists, 
a whole pale swarm of prayer cut into rock until the pillar seemed less carved than murmuring under its own surface. 
Then Light, my side, opposite the gates, too bright, pleased with itself, catching every lamp in polished cuts and little apertures, 
throwing back clever fragments that made me feel chosen because I was fifteen and easily bribed by brightness. 
Then Justice, pale, narrow in spirit though not in size, wrapped in carved chains that crossed and tightened and crossed again, no loose link anywhere, 
nothing generous, nothing accidental, the sort of pillar one did not want to lean against even if one’s feet had begun to burn. \par

%3- The pendulum has two faces. And it turns around itself.

From the upper dark, far above that crowded ring of pillars and stone ribs, the Thread descended, thin as if it had no right to hold what it held, 
a single vertical line dropped through the whole illness of the Hall, past the star, past the capitals, past the greasy lower air of robes and 
lamps and human breath, down to the Pendulum’s disk. 
It was not a ball, whatever foolish picture we made of it when we were small. 
It was a thick round plate, black-rimmed, two-faced, wider than a man was tall, solemn in the idiotic way very heavy things are solemn, because they need no anger to be dangerous. 
One face bore God, hard-cut, old, dry-eyed, less a comfort than a verdict. 
The other was darker, the Beast-face, though we were not encouraged to say it with too much appetite, because children love terrible names and 
adults dislike hearing their own fears enjoyed. 
The disk turned its faces upon the Hall as it crossed, God offered to one side, Beast shown to another, then edge, then thickness, then the other face again, 
a slow exchange of blessings and threats performed by metal and weight.\par

It barely leaned, if one looked all the way up the Thread. 
That was the obscenity of the scale. 
Up there the angle was almost nothing, a small indecency away from the vertical, but below, at the end of one hundred and twenty-five metres of height, 
that little lean became a true crossing, five metres and more from the centre before the disk came back through the middle and went to the other side. 
The empty circle on the floor was not made empty by doctrine first; it was made empty by fear, by calculation, by the common sense of ribs and skulls. 
The barrier around it was low, ugly, badly beloved by every parent in the Hall, and beyond the disk’s own path it kept a margin, so that the forbidden 
floor spread almost eighteen metres across, a round vacancy scrubbed cleaner than the places where people were allowed to stand. 
And the road of the disk did not even stay put. Swing after swing it seemed the same, a back and forth, a return, a patient menace; 
but slowly, too slowly for a bored child and not slowly enough for those who measured such things, the line of travel turned across the floor. 
After two days and a little more, it had made its whole round and come back to itself, as if the Hall had been chewing one invisible circle all that time.\par

% 4-The Hall reveals itself as machinery: hidden water, five Kahir vaults, gravure, and sacred omission

That slow turning gave the two faces their little cruelty. 
No side of the Hall could keep God comfortably before it. 
No side could be spared the Beast forever. 
A family might stand under Balance and receive the dry stare of God at one ceremony, then, later, on another day or another hour, 
find the darker face travelling toward them with the same calm entitlement. 
We children knew this without knowing what to do with it. 
We made secret counts, stupid counts, pious counts, trying to discover whether our side had been favoured, 
whether Light had received the better face longer than Shield, whether Faith had been punished, whether Justice had been made to look at what it deserved. 
The Pendulum ruined all such arithmetic. 
It crossed, returned, turned, and came again, distributing its God and its Beast with the nasty fairness of a machine pretending to be sacred.\par

Around that cleared middle the rest of us had been packed into the Hall’s permitted flesh: Candidates in grey, families in their proper sectors, 
attendants with bowls and cloths, wardens dark and rigid near the gates, old people sitting where age had purchased the right not to pretend at strength, 
children scrubbed too clean and already dirty again at the cuffs. 
I stood on Light’s side, beneath the bright pillar, opposite the closed Mountain gates and the black mass of the Genesis gravure, 
trying to hold my shoulders as if the whole Hall had been assembled to judge their angle. 
The robe scratched my throat. 
The lamps made the air taste of oil. 
Heat gathered under the cloth and stayed there, loyal as a punishment. 
Above me the hollow rose into its dark throat; before me the disk crossed and returned; beneath me my feet began their own small, exact, deeply untheological suffering. 
The ceremony had not even begun.\par

% 5- Ceremony is about to begin, moving or making sound is frowned upon.

Beside me, a boy whose name I have either forgotten or punished out of memory kept moving his lips, not loudly, not enough for correction, 
just enough to wet the words and make himself feel accompanied by them. 

\begin{scripture}
Beyond the forest thou shalt not wander.
\end{scripture}

He had said it twice already. 
Perhaps three times. 
Children repeat scripture the way frightened fingers worry a scab, not always from belief, often because the sentence is there, 
hard and ready-made, better than whatever naked little panic might otherwise come out of the mouth. 
I did not turn toward him. 
Turning in the Hall was never simply turning. 
It made cloth move, made sandals complain, made some adult eye discover that your body had not yet been entirely surrendered to the ceremony. 
So I stayed facing forward, toward the disk, toward the gate-side gravure opposite us, toward everything I did not understand and was expected to honour in the meantime. \par

Then the first chime came from above Faith, or seemed to. 
It did not ring like a little bell in a hand. 
It entered the Hall already enlarged, fattened by stone, thinned again by height, one clear note scraped clean of mercy and sent upward into the hollow 
until the whole almost-cone answered it in layers. 
The sound climbed past the pillars, struck the dark where the Thread vanished, came back along the walls, slid over the annular ring, 
trembled in the carved waves of Faith’s pillar, died in the chains of Justice, and left the lamps hissing louder than before. 
People lowered their heads. 
Robes settled. 
A child swallowed. 
Somewhere near Shield, leather creaked once and then became ashamed of itself. 
The Pendulum crossed again, indifferent to the note, indifferent to us, keeping its own old arithmetic while the human rite finally began to arrange itself beneath it. \par

%6 - Ceremony begins with a procession song.

The Faith chorus started singing, it was going to be a long and tiring ceremony.

\begin{scripture}
Beneath the Mountain we gather,
Beneath the Mountain we sing;
Before the stone and the water,
Before the hidden spring.
Let every footfall be measured,
Let every breath be whole;
Let no one hand hold the whole.

Children of shadow and shelter,
Children of water and stone,
Called to the charge of the Keepers,
Never to stand alone.
Balance the fruit and the cistern,
Faith keep the word in the soul;
Let no one hand hold the whole.

Shield at the gate of the Mountain,
Light where the dark is made known,
Justice between blood and answer,
Counting what each has sown.
Faith lift the prayer through the hollow,
Balance remember the bowl;
Let no one hand hold the whole.

God who has hidden the waters,
God who has covered our head,
Keep us from thirst and from wandering,
Keep us from open dread.
Chain what is bound in the darkness,
Guard what the deep stones hold;
Let no one hand hold the whole.

Now in this solemn procession,
Now as the old names are worn,
One charge is given, one ended,
One cowl lifted, one borne.
Let what is witnessed be faithful,
Let what is spoken be whole;
Let no one hand hold the whole.
\end{scripture}

% 7- Ka-Mereth, Keeper of Faith enters she's close to death.

The last line did not end so much as sink into the Hall, swallowed by height, chewed by stone, returned in little scraps from the pillars until even the words 
seemed to have been handled too much, warmed by too many mouths, made damp with obedience. 
The singers did not look pleased with their own holiness. 
Faith never allowed that. They stood with their throats open and their faces emptied, letting the hymn pass through them as if a person were only a temporary 
inconvenience in the work of sound. 
The bowls in the attendants’ hands trembled faintly. 
The water inside them caught the lamps in yellow skins, broke them, stitched them again, and I remember thinking, with the small stupidity of the body, 
that if one of those bowls spilled, the whole Hall would know it before God did. \par

Then Faith entered properly, not by appearing, for they had already been heard, but by advancing around the Pendulum Circle in their slow white line, 
barefoot, careful, solemn with that practical solemnity of people who have rehearsed danger until it looks like grace. They did not cross the circle. 
No one crossed the circle. 
They curved around it, keeping just beyond the low barrier, carrying water, cloth, a narrow cowl folded over two hands, and behind them came the old 
Keeper of Faith in a low chair of dark wood. 
Ka-Mereth, Keeper of Faith, had not been seen for months, and the Hall made one of those little animal movements crowds make when they have promised 
to be still and are suddenly given something dying to look at. She was not dead. That was the trouble. Death would have been cleaner. 
She had been reduced instead,  narrowed, dried, eaten from within by age or prayer or office, with her white cowl slipping from a skull too small for it, 
her mouth half open, her hands lying on her lap as if they had been placed there by someone who no longer expected them to help. \par

% 8-Va-Elun appears, she is wounded.

Behind the chair walked Va-Elun, Veil of Faith, the one who would not remain Va-Elun for long, and she made the old woman look more terrible by not yet being old. 
Faith had dressed her in white, of course, because Faith believed in white the way sickrooms believe in it, as if paleness could bully the body into purity. 
But over the robe she wore the pilgrim’s mantle, rough, undyed, mended in darker patches, the kind of garment that had been dragged through stone and sweat 
and prayer until it no longer belonged to fabric entirely. 
At her throat, where the collar opened, two raw lines crossed the skin. 
Not ornaments. 
Not symbols yet. 
Wounds first, red and ugly and badly healed, which meant the Hall could make them sacred only after everyone had been forced to see that they had hurt.\par

I stared at them. 
I was fifteen; I stared at wounds, at women, at authority, at anything adults had not yet taught me how to look at without giving myself away. 
The scars seemed indecent in that white procession, more indecent than nakedness would have been, because they were not offered and yet had 
been put where all of us could receive them. 
Va-Elun kept her face lowered, emptied like the singers’ faces, but the throat betrayed her, the little pull of the skin when she swallowed, 
the life still working there under Faith’s cowl and mantle and all that washable holiness. 
I might have gone on staring until some future version of myself died of shame, but Va-Raedin, Veil of Light, touched two fingers to the back of my wrist, 
lightly, almost kindly, which in Light was often the cruelest available form of correction.\par

% 9-Kailan stares at the wound, and Va-Raedin corrects him into Light’s first lesson

Va-Raedin did not even look at me when he corrected me, and the absence of his gaze deprived my shame of the small consolation of being an event; 
no public rebuke, no lifted eyebrow, no splendid cruelty to resent, only two fingers placed against the back of my wrist with the hygienic delicacy 
of a man removing a stain from glass. 
I lowered my eyes to the stone, to the hem of my robe, to the bright stupid edge of my sandal, and for a moment the whole Hall shrank to that: 
leather biting the foot, grey cloth, bloodless knuckles, the disgrace of having been caught feeding on another person’s pain as if ceremony had served it to me.\par

The procession reached the Faith side of the circle and began its careful arrangement, all that holy fuss of bodies and objects finding their assigned places: 
the bowls set down on low stones, the folded cowl held higher than it deserved, attendants bending, straightening, turning with their thin white sleeves 
whispering against themselves. 
Ka-Mereth’s chair was carried not to the centre, never to the centre, but to the edge of the pillar-ring, where the Pendulum’s road could pass before her without claiming her. 
The old Keeper’s head slouched once, then righted itself with a horrible little effort, and the movement went through the Hall more surely than a shout. 
It told us she was still inside the body. 
It told us the ceremony would have to use her before it released her.\par

% 07-The five named waters are placed, and Faith prepares the transfer

Va-Elun knelt before her, and the kneeling altered her more violently than standing had. 
Upright, she had still been a woman approaching office; kneeling, she became material for it, neck offered, hands open, mantle rough against white cloth, 
scars pulled pale where the skin bent. 
Faith loved these postures in which the body appeared to consent to its own use. 
The bowls had been set around them, five small mouths of water waiting on stone, each trembling with the last echo of the hymn. 
In their surfaces the lamps spoiled themselves, lengthened, broke, came together again in yellow streaks, and the old Keeper, 
the new one, the attendants, the pillars, even the black rim of the passing disk were all reduced there to small untrustworthy stains.\par

For a while nothing happened, or what happened was so slight that the Hall had to become greedy for it. 
Ka-Mereth breathed. 
That was the event. 
Her mouth opened, closed, opened again with the dry patience of a cracked vessel being asked to pour. 
An attendant bent toward her, offering perhaps a word, perhaps merely the odour of a living face, but the old woman lifted two fingers and the attendant withdrew at once, 
chastened by that tiny remainder of command. 
Then Ka-Mereth raised her head. 
The movement was dreadful because it took so long. 
One could see the effort travel through her: throat, jaw, eyelids, the whole poor machinery of survival dragged into service by office, 
until at last her eyes settled on Va-Elun and the Hall leaned inward without moving.\par

% 08-Ka-Mereth’s failing voice begins the Faith catechism: gift, refusal, Temple, and forbidden silence

Then Ka-Mereth spoke, or rather allowed a thread of voice to escape the ruin of her mouth, thin, scraped, unexpectedly clean, 
the sort of voice that made one ashamed of having expected weakness to sound weak. \par

"What did God give the hidden people?"\par

"Stone, water, and shadow," the Hall answered, and the answer did not burst upward like faith, no, it came out thickly, obediently, 
from mouths already trained to the shape of it: old mouths with saliva gathered in the corners, children’s mouths still pink from scrubbing, 
bored mouths, devout mouths, frightened mouths, the sour and the pure all yoked for a moment to the same formula. The words climbed the tiers, 
struck the walls, returned diminished, and for the length of that answer the Mountain performed its old fraud of unity, making us sound like 
one body at precisely the moment we were reciting the gifts by which that body had been cut into charges, sectors, offices, permissions, 
little sacred amputations of the whole.\par

Ka-Mereth’s eyelids trembled. 
Her mouth worked again, dryly, as if the next question had to be fetched from behind the teeth. 

"What did God refuse the hidden people?"\par

"Sky, wandering, and the open world." came the answer, and something passed through the Hall then, not disobedience, 
never that, but the small inward illness produced by naming a thing one has been ordered to call a danger before one 
has even been allowed to desire it properly. 

I had never seen the open world then, not truly; not Switzerland with its insulting windows, its streets that did not count 
a body before releasing it, its taps that gave water without witness, its churches standing there in public like mistakes 
no Shield had corrected. But even as a boy, before any of that, the three words had their taste. 
Sky. 
Wandering. 
Open world. 
They were supposed to fall from the mouth like refusals. 
Instead they shone there briefly, poisonous and blue, three luxuries dressed in the rags of prohibition, three hungers made pious by being denied.\par

The old Keeper breathed again, and the Hall waited for the part it knew best. \par

"What is kept beneath the Mountain?"\par

"The Temple."\par

"What is kept beneath the Temple? \par

No one answered. 
The silence arrived complete, not as an absence but as a substance poured instantly into every mouth. 
It filled us better than song. 
Even the children knew how to hold it, lips sealed, tongue stilled, breath made careful behind the teeth, because this was one of the first arts taught in the Mountain: 
not reverence, not courage, not understanding, but closure. 
Some questions were asked only to harden the refusal of speech, to train silence until it became muscular, hereditary, almost proud of its own deformity. 
Ka-Mereth sat inside that silence like a relic that had not yet been permitted to finish dying, and Va-Elun knelt before her, receiving not an answer but 
the old black hollow where the answer had been buried.\par

% 09-Va-Elun kneels before Ka-Mereth.

The silence held until Ka-Mereth broke it herself, not by answering, which would have been a vulgarity, but by lowering one hand toward the nearest bowl. 
An attendant tried to help her and was repelled by the smallest twitch of the wrist, a movement so dry, so insufficient, that it seemed less like refusal 
than the memory of refusal surviving in bone. 
Her fingers touched the water. They did not enter it grandly. They spoiled its surface, that was all, 
wrinkling the yellow reflections of the lamps, making them shudder into little broken organs of light. 
Then she lifted the wet fingers above Va-Elun’s bowed head, and the drops fell with a miserable precision, one on the hair, one on the cowl, 
one on the raw red crossing at the throat, where water found the wound and made it shine. \par

Va-Elun did not move. That was her triumph, I suppose, though at fifteen I understood triumph mostly as noise, speed, victory, anything with the decency to announce itself. 
This was another sort: the triumph of not flinching while an old office wetted your scars in front of everyone and made your private pain communal property. 
The Hall watched with that purified greed ceremonies permit, the greed of good people allowed to stare because staring had been renamed witness. 
Around the circle the Pendulum crossed again, God-face, edge, Beast-face, indifferent to the little sprinkling below, 
and the Beast seemed for a breath to look not at Va-Elun but at all of us, at our washed mouths, our sealed questions, our obedient hunger for what would happen next. \par

%10- Q&A about the chamber of the last prayer.

Ka-Mereth moved one finger.
That was enough. 
The attendants around her stiffened as if a thread had been pulled through all their spines, and the nearest bent so low his hood almost touched the arm of her chair. 
She whispered into the cloth beside his ear. 
We heard nothing, which made the words enormous. 
The Hall, always hungry for what it was not permitted to swallow, leaned inward by a thousand invisible muscles. 
Then the attendant straightened and lent her his healthier voice. \par

"Va-Elun, Veil of Faith." Va-Elun was already kneeling. She had been there so long that I had almost ceased to see her, 
a white shape folded before Ka-Mereth’s chair, low as penitence, still as discarded linen after some difficult washing of the soul. 
Her robe had been cleaned for the Hall, but not redeemed. 
Dust remained in the seams. 
The hem was grey with the powdered stone of lower passages. 
At one cuff, old water had dried into a brownish stiffness, the sort of stain adults pretend not to notice when it has come from a holy place. 
Her head was bowed. 
Her hands lay open on her knees, palms upward, empty, as if emptiness itself had been demanded below and she had returned carrying it with perfect obedience.
The attendant spoke again. \par

"Where did the path end?" Va-Elun did not lift her head. \par
"Where asking ends." \par

A small unease moved through the Hall, not sound, not disorder, only that collective inward rot by which a crowd understands that a sentence has passed above 
it and refused to be understood. 
The words seemed to have come up from the labyrinth still damp, still cold, with a little blackness clinging to their edges. \par

"What did the dark give you?" \par
"No answer." \par
"What did you give the dark?" \par
"No demand." \par

Ka-Mereth’s eyelids lifted.
Her eyes had been almost extinguished, folded back into the private twilight of a body near its end; now they returned, narrow and cruel with recognition, 
as if some buried mechanism in her had heard the correct pressure and begun again. 
The attendant bent once more. 
Again the tiny whisper. 
Again his borrowed voice carrying what her own body could no longer afford. \par

"What stands in the Chamber of the Last Prayer" \par
"What cannot be carried." \par
"What returns from it?" \par
"What was never mine." \par

No one moved. 
Even the children understood that this was not a prayer made for us. 
It belonged to some older traffic between office and darkness, some underground catechism kept where ordinary words went to moulder. 
Ka-Mereth stared down at Va-Elun. 
The Pendulum crossed and returned. 
The Hall waited, full of its own breath, its oil, its old dust, its obedient ignorance.
Then the dying woman made the same small motion with her finger.
The attendant turned toward the lower passage. \par

"Bring what returns," he said. \par

%11- The bowl of the last prayer
The Bowl of the Last Prayer had not been seen in a lifetime. \par

That was why the murmur escaped before discipline could throttle it. 
Everyone knew the name: children as they knew sealed doors and punishments kept for later; old women touching thumb to finger, 
one, two, three, four, five; men who affected unbelief lowering their voices all the same. 
The bowl belonged below, in the Chamber of the Last Prayer, beyond the blind turns of the labyrinth, where no lamp was carried and voices came back altered. 
It was brought out only when a Veil of Faith had to prove, before us all, that the last path had been walked, the chamber found, and the holy thing returned 
from its proper darkness. \par

They uncovered it last, and the Hall leaned toward it like a crowd around sickness. 
The bowl sat in a black cradle, elliptical, pale, smooth only where the water would touch it. 
Outside, the alabaster had been tormented into sanctity: five bands of scripture, hooded figures no larger than fingernails, 
gates, roots, chains, waves, jars, lamps, eyes, hands, little mountains, little mouths, all cut into the stone until no innocent surface remained. 
Gold slept in some grooves. 
Ash clotted others. 
Old oil had browned the lower carvings into a greasy amber. 
In places the alabaster was so thin that lamplight entered and curdled there, like old milk under skin.\par

Inside, the bowl had been spared. 
After all that outer delirium, the interior opened into a long clean hollow, polished to a milky silence, bare of word or figure. 
Faith had decorated everything except the thing that mattered. 
The outside belonged to worship. 
The inside belonged to water. 
And there, rising from the pale floor near either end, stood two little columns of dark metal, no taller than a finger joint: the Ear and the Mouth. 
Everyone knew their names. \par

They looked ridiculous at first, like two pillars from an insect chapel, and then less ridiculous, because everything in the Mountain became more obscene when reduced. 
The metal was almost black, old singing bronze, drinking the lamps instead of returning them. The chamber water lay shallow around their feet. 
The Ear had been darkened by touch. The Mouth looked cleaner, meaner, a sealed little throat waiting to deny it had ever spoken.\par

Ka-Mereth was carried close enough to look inside. 
One attendant presented a narrow rod of the same black metal. 
Another steadied the cradle. 
Behind them, the new Veil of Faith stood under her hood, too pale, too emptied, having returned from the labyrinth with the bowl and with whatever the 
labyrinth had taken from her in payment. 
She did not look triumphant. 
She looked like a man who had learned that some doors open only to make return less believable.\par

Ka-Mereth took the rod. 
Her hand shook so violently that two attendants moved to help, but she flicked them away with a dry remnant of a gesture, and they obeyed. 
She lowered the rod toward the Ear. 
The Pendulum crossed the centre of the Hall with its enormous mute assurance. 
Then her hand dropped the last finger’s width.\par

The sound was almost nothing: a small dry tick of metal on metal, miserable after so much expectation. 
A few children shifted. 
I remember the shame of having expected thunder.
Then the water moved.
The first rings left the Ear, widened, touched the clean alabaster curve, and returned wrongly. 
They ran along the pale wall, bent inward, crossed themselves, thinned, gathered, and for one instant made near the Mouth a shape I could 
not name then and cannot honestly name now without stealing from later knowledge. 
It was not a picture. It was not writing. It was a little treason of water against the common law of the eye.\par

Then the Mouth answered. 
Not loudly. 
Never loudly. 
A thin note came from it, so high and narrow that it seemed less heard than discovered inside the teeth. 
It lasted only a breath. 
Then it died, and the water smoothed itself again with the hateful innocence of things that have proved a point and refuse explanation. 
Ka-Mereth’s eyes, almost extinct a moment before, sharpened with private violence. 
She lowered the rod, still staring at the Mouth. \par

"The chamber answers," she said. \par

The words passed through the Hall with more force than a shout. 
The Veil of Faith bowed his head. 
The bowl had not been admired. 
It had not been blessed. 
It had been tested, and something in it, something below it, something behind the engraved clutter of our ignorance, had replied.\par

%12- Va-Elun is named Ka-Elun, Keeper of Faith.

For a moment after Ka-Mereth spoke, no one seemed to know what to do with their body.
The Hall had received the words and could not digest them. 
"The chamber answers." 
They remained in the air, not as comfort, not as triumph, but as one of those old verdicts that make obedience easier by making thought unpleasant. 
Va-Elun stayed kneeling. 
The Mouth of the bowl had gone silent. 
The water had recovered its innocence. 
The Pendulum crossed and returned, carrying its huge indifference over our small success, and the lamps hissed on, vulgar as ever, because lamps have no respect for revelation.\par

Ka-Mereth lifted the rod a second time.
Not toward the Ear. Toward Va-Elun.
The gesture was so slight that another woman’s hand might have made nothing of it, but from Ka-Mereth it became command, accusation, inheritance. 
Va-Elun lowered herself further, until her forehead nearly touched the stone. 
The attendant bent to the dying Keeper. 
Again the whisper went into cloth. 
Again the borrowed voice came out larger, cleaner, more shameful for being strong where hers was failing.

"What answered below has answered above."
"What was hidden has been returned."
"What was never yours is now entrusted to you."

Then Ka-Mereth herself tried to speak. 
The first sound broke in her throat, dry and ugly, and the Hall suffered it with her. 
She gathered what remained. 
Her mouth opened again. 
This time the words came out thin, almost spoiled by breath, but they were hers, and because they were hers they struck harder than the attendant’s polished voice.

"Rise as Ka-Elun, Keeper of Faith." \par

Ka-Elun rose.
Slowly, badly, not like a figure in a carving, not like the clean ascent the songs pretended. 
Her knees had stiffened. 
Her robe dragged at the stone. 
Dust fell from one fold in a little grey sigh. 
When she stood, she seemed taller only because the kneeling had ended, not because any glory had entered her. 
An attendant removed the Veil’s white hood from her head. 
Another brought the Keeper’s cowl, heavy with old embroidery, yellowed at the edges, weighted by thread, smoke, hands, years. 
It was lowered over her shoulders like a burden rescued from a corpse. \par

Mereth watched it settle.
Her face did not soften. Faith had no use for softness at such moments. 
The old woman looked at the new Keeper as if measuring whether the name had fitted badly enough to be true. 
Then her hand fell back into her lap, empty now, more terrible than when it had held the rod.
Ka-Elun bowed first to Mereth, then to the bowl, then to the dark lower passage from which it had come.
Only after that did the Hall breathe.
Not all at once. 
Breath returned by sectors, by families, by frightened children discovering that their mouths still belonged to them. 
Faith had changed hands before us. 
The chamber had answered. 
The relic had sung. 
A woman had risen with another name on her, and the old Keeper remained in her chair, diminished but not released, 
like a candle that has lit another and resents the flame for surviving elsewhere.\par


% 11-Kailan sees Va-Sheva: desire, rank, humiliation, and Shield authority fused together

I looked toward Shield then, not because Shield could answer anything, but because my eyes had already learned 
the path to Va-Sheva before my mind could invent a better excuse, and I found her at once beside the northern 
arch, four steps below her Keeper and above the armed wardens, dark green-black leather over the sleeveless 
robe, hair bound high, short blade at the hip, longer one across the back, standing there with the awful 
efficiency of someone shaped before others had finished becoming.\par

She was nineteen, perhaps twenty, old enough that the women beside her treated her as one of them, old enough
 that the boys near me watched her with that careful hopelessness boys reserve for things already claimed by 
 adulthood, and I watched too, though I would have denied it with all the stupidity available to me, because 
 to look at her was to feel my own incompletion, my narrow shoulders, my eager face, my whole unfinished self 
 standing in a robe that still belonged more to ambition than to office.\par

When she looked in my direction, her eyes passed over me, paused because recognition required it, and softened 
by a fraction, not enough, never enough, and that was the cruelty of her kindness, that she never gave more 
than truth allowed, never pretended the gap between us had closed, never smiled at the boy I wanted to be 
when she could still see the child who had followed her through service tunnels and cried once over cut lamp 
glass before denying it until she laughed.\par

Wanting her then was not yet what men call wanting when they wish to make themselves sound simple, it was 
humiliation, rank, distance, hunger for approval, the pain of being measured by someone whose measure mattered,
 and perhaps, yes, desire too, though desire was the least honest word for it, too clean, too proud, 
 too ready to flatter itself, while what I felt was closer to trespass upon a threshold I had not earned 
 the right to approach.\par

She looked away before I did.\par
That was mercy, perhaps.\par
Or indifference.\par

At fifteen, I did not yet possess reliable instruments for distinguishing the two, and it is one of exile’s 
small cruelties that later, when you finally learn the difference, the people whose glances taught you are 
no longer standing where memory can test itself against them.\par

% 12-The five Keepers witness Faith’s transfer: Balance measures, Shield seals, Justice records, Light illuminates

On the dais, the rite continued, because rites do not care what they interrupt inside boys, and the five 
Keepers, or those who served in their place where age or office required, came forward to witness Faith’s 
transfer, one by one, each Kahir presenting not only itself but its suspicion of the others, because that 
was what we called balance before I knew enough words to call it containment.\par

Ka-Leth of Balance came first, broad, heavy, grey at the temples, her hands stained faintly green from 
ledger dye and plant work, and two attendants followed with a narrow measuring rod and a copper vessel 
marked by rings, so that even at the moment when prayer was supposed to rise beyond measure, Balance 
stood there to remind the Hall that a blessing which emptied the cistern would still leave thirsty mouths 
behind it.\par

She poured a thread of water from her bowl into the central basin, and she waited until it touched the 
first line, not above, not below, because mercy in Balance was not softness but accuracy, and only after 
the line had been met did she speak.\par

“Let prayer not empty the cistern,” she said.\par

Ka-Dhavar of Shield came next, tall, unsmiling, with scars crossing one hand, and he did not pour, of 
course he did not pour, Shield gave nothing fluid to a rite if iron could speak more truthfully, and he 
placed beside the basin a sealed ring that looked too plain to be sacred and too heavy to be symbolic, 
while at the northern arch three wardens changed stance together, one inward, one outward, one upward 
toward the galleries, and no one commented because everyone understood that Shield’s prayers were made 
first with bodies.\par

“Let no threshold open because a holy tongue grows soft,” Ka-Dhavar said.\par

Then Justice, pale-robed, cold, without ornament except the strip of black cord at the wrist, Ka-Orun 
carrying a tablet of polished stone and reading the transfer into the record before it had fully happened, 
which unsettled me then and unsettles me still, because Justice had the indecency to write events while 
they were becoming themselves, catching the breath before grief could improve it, placing sequence around 
the living like a frame around an animal not yet dead.\par

“Let what is witnessed be kept,” Ka-Orun said.\par

“Let what is kept be answerable.”\par

Then Light came forward, and my foolish heart leaned toward him before my body did, Ka-Syphiron, 
Keeper of Light, old but not ruined like Ka-Mereth, white-gold robe catching the lamps and returning 
them with such irritating grace that even age seemed to have become an angle under his command, a beautiful 
old man in the way polished bone is beautiful, severe, pale, almost cruel in its clarity, and behind 
him Va-Raedin stood still as a blade laid flat.\par

Ka-Syphiron lifted a lamp before Va-Elun’s bowed head, and the flame inside it burned blue at the root, 
which I saw and wanted immediately to admire, wanted to believe in, wanted to be seen understanding before 
I understood anything at all, and though he did not speak to me, Va-Raedin’s earlier correction returned 
so sharply that I felt it at the back of my wrist again: do not admire the flame, ask where the air enters.\par

“Let the path be seen only by the one appointed to walk it,” Ka-Syphiron said.\par

Faith’s attendants removed the old Keeper’s cowl, and for a breath Ka-Mereth’s bare head was shown to the 
people, thin hair, white scalp, the mortal indignity beneath sacred office, and this too the Mountain 
required, not out of cruelty exactly but out of that worse instinct by which institutions make the body 
confess its weakness so the office may appear merciful for surviving it.\par

% 13-Office edits the body: Va-Elun disappears into the cowl and rises as Ka-Elun

Then they placed the cowl over Va-Elun, and it was too large for her, falling around her shoulders, covering 
the torn mantle, the patched sleeve, the marks at her throat, swallowing the evidence of the path beneath 
the cloth of arrival, and I saw then without knowing how to say it that office does not merely crown a body, 
it edits it.\par

Ka-Mereth lifted one hand.\par
“Rise now as Ka-Elun, Keeper of Faith.”\par

The words were almost nothing, less than a command, less even than a final breath properly spent, but the 
Hall made them enough, caught them, carried them, distributed them, and when Va-Elun stood she seemed at 
first to disappear inside the office, as if the cowl had not been placed on her but over the space where 
a person had been.\par

Then Ka-Elun spoke.\par

Her voice was stronger than the old Keeper’s, but quieter than I expected, and that quietness did more 
work than force would have done, because she did not climb toward the Hall, she made the Hall come toward 
her, made us lean inward, made hearing into obedience, which is one of Faith’s oldest arts and one of 
the reasons it survived so long beneath stone.\par

“The Mountain keeps what the open world destroys,” she said.\par
“The Oasis remains because it is hidden.\par
The Temple remains because hands are divided.\par
The people remain because memory does not wander.”\par
“Memory does not wander,” the Hall answered.\par

I said it too, with the others, though even then some small disobedient part of me wondered what 
kind of memory must be guarded so fiercely from movement, and later, in rooms where Swiss clocks cut the 
day into bright indifferent pieces, I would learn that memory wanders whether one permits it or not, 
that exile is memory escaping its assigned chamber and carrying contraband into every clean room it enters.\par

% 14-Faith states hidden survival, divided hands, darkness as inheritance, and memory under command

Ka-Elun continued, and now the words were from The Books of The Covenant, though not one of the 
passages usually given to children.\par

\begin{scripture}
“And God said, Fear not the darkness which I have given thee; for the darkness shall be thy veil, and the stone shall be thy wall, and the deep waters shall be thy life.”
\end{scripture}

I knew the line, everyone knew the line, we had heard it under sleep, over water, beside burial, 
in the mouths of old women and teachers and Faith attendants washing the faces of children, but 
spoken there by the new Keeper with the old one still breathing behind her, it altered its weight, 
and darkness ceased for me, briefly, to be the absence children feared and became something given, 
distributed, inherited, perhaps even administered.\par

“Mother of Faith,” Ka-Elun said, though no blood joined them, “your charge is complete.”\par

Ka-Mereth’s lips moved, but no sound came, and the attendant bent close to receive whatever 
remained of the old voice, then stood and spoke in her place.\par

“She goes into prayer.”\par

The Hall lowered its heads, and the old Keeper was carried away through the Faith arch toward 
chambers I had never seen, would not see, perhaps was never meant even to imagine, quiet cells, 
warm walls, prayer benches, attendants trained to wash bodies that had once commanded the Mountain, 
rooms for those whose authority had outlived strength and whose strength had outlived usefulness.\par

Children were told such retirement was holy.\par

I thought it looked like being sealed slowly while still alive, and I hated myself for thinking 
it, which did not stop the thought, which is how I first learned that shame is not repentance, 
only the body’s alarm at having produced knowledge before permission.\par

% 15-The ceremony does not end: Kailan’s own Ascension begins after Sa-Tavan’s Justice naming

Then the ceremony changed.\par

It did not end, and that was the first warning, because the great Ascension had been completed, 
Faith had received its new Keeper, the old one had been carried into prayer, the Hall had answered, 
the words had done their work, but the Kahir remained in place, the bowls stayed on the stones, 
the bronze grille still whispered behind Balance, Shield did not release the exits, and somewhere 
behind my shoulder Va-Raedin’s hand withdrew not in disapproval but in attention.\par

“Now,” he said.\par

My mouth went dry.\par

I had known this was coming, of course I had known, my mother had washed my robe herself that 
morning despite the attendants Light had sent, my father had held my wrist before leaving our 
chamber and told me to stand straight when spoken to, my younger cousins had stared at me with 
the proud stupidity of children adjacent to a story, but knowing is a poor instrument when the 
thing known begins moving toward the body.\par

The names of the young were not the heart of the day, and that made it worse, because a private 
elevation might have belonged to me, while this was a smaller rite attached to a greater continuity, 
a young growth bound onto old root, and there is humiliation in being told you are important only 
as evidence that something larger than you has not yet failed.\par

Sa-Tavan went first, presented before the Justice alcove, and I watched him answer with that 
careful, almost irritating precision of his, a boy already made angular by the office that had 
chosen him, and when Ka-Orun asked whether he would keep peace without mistaking silence for peace, 
Tavan paused, only once, barely enough for breath to notice, and said, “I will try.”\par

A faint movement passed through Justice, not displeasure, not approval, something colder and more 
pleased than both, because Justice distrusts certainty in the young, and perhaps in the old too, 
though the old are better at disguising certainty as method, and Ka-Orun tied a narrow black cord 
around Tavan’s wrist.\par

“Rise now as Sa-Tavan, Candidate of Justice.”\par

% 16-Kailan descends to the eastern stone, becomes visible to himself, and faces the unrehearsed question

Then it was my turn.\par

I descended into the Hall and the Hall became larger with every step, the faces rising in tiers, 
workers with damp hair, Balance clerks, Shield wardens, Faith singers, children on hips, old men 
with polished canes, my parents somewhere among them, and all of it blurred at the edges except 
the eastern stone ahead and the knowledge that if I stumbled I would remember it until death.\par

I passed below the Shield arch, and Va-Sheva was close enough for me to become aware of everything 
wrong with myself at once, the rebellious hair on the left side of my head, the dust on the edge 
of my sandals, the eagerness my face produced when I forgot to discipline it, the whole offensive 
evidence that I was still a boy moving through a rite that wanted to call me something sharper.\par

She did not smile, because it would have been improper, but her eyes found mine for one instant, 
and that was enough to ruin my balance more effectively than any uneven stone, because there are 
glances that do not give hope but make hope behave indecently, standing up inside you at the 
exact moment you need it to kneel.\par

Va-Raedin led me to the eastern stone.\par

Ka-Syphiron watched me approach with those pale eyes of his, unpleasantly clear, eyes that seemed 
not to look at a thing but through the excuses a thing had prepared in case it was seen.\par

“Kailan,” he said.\par

I knelt.\par

The stone was warm under my knees, too warm, someone had opened the vent beneath the eastern 
dais more than ceremony required, and I remember this because fear attaches itself to details, 
to heat, to breath, to the ridge of a sandal, to anything smaller than the question waiting above it.\par

“Do you seek Light?”\par

“I seek Light.”\par

“Why?”\par

That question was not in the rehearsal.\par

The Hall changed around it, not loudly, not visibly, but by that tightening of attention every 
gathered body knows before the mind admits it, and I felt Va-Raedin behind me, felt my parents above, 
felt Sa-Tavan newly corded, felt Va-Sheva somewhere to my left, older, sharper, impossible, and I 
understood that the safe answers had all been removed from the table before I knew we had left the script.\par

Because I am worthy was too proud, because I want to know too naked, because darkness is not 
obedience too dangerous, because I have read everything a Candidate is permitted to read and it has 
taught me mostly the shape of what I am not permitted to read absolutely suicidal, and yet the words 
were moving in me faster than caution, arranging themselves before fear could inspect them.\par

% 17-Kailan gives the dangerous answer: hidden things still have shape and may be found

“Because what is hidden still has shape,” I said.\par
“And if it has shape, it may be found.”\par
The Hall went silent.\par
Not the silence of reverence.\par
Not even the silence of shock.\par %And possibly, given this is a very Faith principle, Faith makes a weird notice. But we don;t understand what's this about

A worse silence, an administrative silence, a silence already making copies of itself for later rooms, 
because everyone present understood that something had been said which would not remain inside the ceremony, 
something small enough to deny and large enough to record.\par

Ka-Syphiron’s expression did not change, which frightened me more than disapproval would have done, and 
behind me Va-Raedin’s presence altered by a degree, his hand no longer resting where it had been, not 
withdrawing from me but attending more closely, as if I had finally become a problem worth the full courtesy 
of his notice.\par

“A Candidate does not find all that has shape,” Ka-Syphiron said.\par
“No, Keeper.”\par
“What does a Candidate learn?”\par
Relief came so quickly it almost made me stupid.\par
The script had returned.\par
“Where light falls.”\par
“And what does a Veil learn?”\par
“Who placed it.”\par
“And what does a Keeper learn?”\par

Here I was meant to say what must remain dark, and I knew the answer, knew it from instruction halls, 
correction cells, whispered contests among boys who wanted to sound old and wounded before life had earned 
them either condition, but after Ka-Elun’s torn mantle, after Ka-Mereth carried away alive into prayer, 
after the Chamber of the Last Prayer had crossed the Hall as a real place rather than a story, the answer 
no longer sounded wise.\par

It sounded like a door closing.\par
“What must remain dark,” I said.\par

% 18-Kailan is cut, blood-bound to Light, named Sa-Kailan, and acknowledged by the Mountain

Ka-Syphiron studied me for another moment, and I have often wondered, since, what exactly he saw there, 
whether he saw only a boy late in obedience, a boy whose mouth had reached the right answer after the 
safe moment had passed, or whether Light possesses, in its old men, instruments more exact than sight, 
capable of measuring the interval between submission and understanding, between doctrine repeated and 
doctrine believed, but he did not correct me, not then, he only lifted his left hand, and Va-Raedin 
placed a blade into it.\par

It was not ornate.\par

That mattered.\par

Light did not decorate the thing that cut.\par

The handle was bone, worn smooth by fingers older than mine, and the edge was bright in a way the lamps did not explain but merely revealed, a narrow fact, waiting.\par

I extended my hand.\par

There are moments the body understands before the mind, and mine understood that it was being 
admitted into doctrine not as witness but as material, the way oil, water, wool, stone, and 
breath were admitted, things to be used because they could be measured, moved, burned, spent, 
or marked.\par

The cut came across the base of my thumb.\par

Quickly.\par

Too quickly for courage.\par

Pain arrived first as surprise, then heat, then a kind of occupation, as if a small foreign kingdom 
had opened under the skin, and blood gathered under the lamplight, dark, thicker than I expected, 
stupidly private, embarrassingly mine, though the whole point of the rite was to make it stop belonging only to me.\par

Ka-Syphiron turned my hand with the care one gives to an object already entered into use, and pressed 
the bleeding thumb against a thin plate of polished stone where five faint circles had been marked, 
so pale I had not seen them until my blood moved toward one of them, and touched only the eastern mark, 
only Light, only the office that had praised my hunger by cutting the part of me that reached.\par

The stone was cold.\par

That is the detail memory kept.\par

Not the words first.\par

The cold.\par

My blood was warm, the stone was cold, and between them the Mountain made a contract it did not ask me to read.\par

“Blood enters service,” Ka-Syphiron said.\par

“Sight enters discipline.\par

Pride enters correction.”\par

I repeated the words because repetition makes terror useful, because a mouth busy with formula has less 
room for panic, because the Mountain knew, better than any physician, that the wounded body will obey a 
sentence if the sentence is given quickly enough.\par

“Blood enters service.\par
Sight enters discipline.\par
Pride enters correction.”\par

The words were different with my blood on stone.\par

All words are.\par

A sentence spoken before blood is instruction.\par

A sentence spoken after blood is law.\par

Ka-Syphiron lifted his hand away from mine, and for one small instant the print remained there, red on the eastern circle, a boy reduced to a mark, a mark raised into office, and the absurdity of it struck me only years later, in a country where men signed papers with pens and believed that made them more reasonable than blood.\par

“Rise now as Sa-Kailan, Candidate of Light.”\par

The name did not fall on me like praise, not as I had rehearsed it in the foolish theatre of childhood, not warm, not radiant, not like the Mountain opening its arms to what I had always secretly been, but like a door closing with me on the correct side, or the wrong one, I could not yet tell, and I rose feeling not larger but entered: entered into records, entered into permissions, entered into correction, entered into eastern rooms whose locks would now recognize me.\par

The Hall did not cheer.\par

Cheering belonged to the stories we told about Outsiders and their kings, noisy people who needed sound to convince themselves that an event had happened.\par

The Mountain acknowledged.\par

That was different.\par

The people struck two fingers against stone, once, twice, five times, and the sound filled the Hall like rain made hard, not applause, never applause, but agreement turned mineral, a thousand small impacts accepting that the old order had made room for one more wound.\par

I stood inside it with my thumb throbbing and my robe suddenly too light for my body.\par

I had thought I would feel transformed.\par

Instead I felt cut, watched, and unfinished.\par

Va-Raedin leaned close enough that only I could hear him.\par

“You answered poorly,” he said.\par

My stomach dropped so cleanly that for one breath I forgot the blood.\par

Then he added, “But not uninterestingly.”\par

That was worse.\par

It was also not entirely bad.\par

Light, I would learn, had a genius for that particular cruelty, the wound that also feeds, the insult that opens a door, the correction one stores like contraband because it is the nearest thing to affection the office permits.\par

I did not yet know what to do with a thing that could diminish and enlarge me in the same instant.\par

Much of my life would later be spent among such things.\par

% 19-Faith closes the rite while Kailan repeats the divided-hand doctrine and already begins to doubt it

The final portion of the rite belonged to Faith again, because Faith closed what others opened, and Ka-Elun stood before the five bowls while the Hall endured the last stretch of its own devotion, heat collected under the vaults, children asleep against adult shoulders, old people grey with fatigue, oil thick in the air, the crowd’s body beginning to remember that holiness too must negotiate with knees, thirst, wool, breath, and the longing for release.\par

\begin{scripture}
“God who hid the waters beneath the Mountain,” Ka-Elun said, “keep us faithful in measure.
God who heard the first blood cry from the dust, keep us faithful in consequence.
God who set the desert before the stranger, keep us faithful at the threshold.
God who gave the Gospel beneath stone, keep us faithful in memory. God who divided light from darkness,
keep us faithful in sight.”
\end{scripture}

\begin{scripture}
“Let no one hand hold the whole.”
\end{scripture}

The Hall answered.\par

“Let no one hand hold the whole.”\par

I said it with them.\par

Of course I did.\par

One says such things when one’s thumb is open, when one’s new name is still raw, when thousands of mouths move around one’s own and the body has not yet distinguished belief from rhythm, but even as I repeated it, some small unlicensed part of me turned the sentence over as if it were a token from the eastern rooms, no one hand, the whole, the division made sacred, the ignorance made duty, and I wondered, not nobly, not bravely, only privately, whether Faith had found truth there or only the most beautiful possible name for fear of anyone seeing enough to object.\par

That was the first disloyalty of the day.\par

Not my answer to Ka-Syphiron.\par

Not the blood.\par

This.\par

The obedient mouth, and behind it the small hand inside the mind testing the lock.\par

Then I looked toward Va-Sheva and wondered what kind of chamber Shield kept.\par

That is how quickly the mind betrays solemnity.\par

Faith had just sealed the rite, Light had cut my hand, the Mountain had spoken my new name, and still I looked for the woman at the northern arch, except she was no longer watching the ceremony, no longer even pretending to belong to its ending, she was looking toward the western exit, and two wardens had come close to her, one speaking without moving his head, the way Shield speaks when words themselves must not become visible.\par

Her expression tightened.\par

Not fear.\par

Readiness.\par

There was almost nothing to see, which is why I remember it so clearly.\par

A hand near the hilt.\par

A glance past the galleries.\par

A small adjustment of stance that changed the space around her.\par

Shield did not require drama to become Shield.\par

It became Shield by withdrawing mercy from uncertainty.\par

She touched the hilt at her side and nodded once, and I noticed it because I was of Light now, yes, that is the answer the new Candidate would have preferred, trained to observe what others missed, but the truer answer is less flattering and has survived every correction since: I noticed because I had always noticed her.\par

% 20-The ceremony dissolves: Va-Sheva shifts into Shield, Va-Raedin gives the first instruction, and Sa-Kailan follows into Light

The ceremony dissolved slowly, which is to say it dissolved according to rank, because even departure in the Mountain had to prove that nothing had become vulnerable merely because the sacred part had ended.\par

Keepers withdrew through their arches.\par

Families waited for their tiers to be released.\par

Faith attendants lifted the bowls with both hands, careful not to spill water already named by office.\par

Balance checked the basin levels before permitting removal.\par

Justice sealed one tablet before opening the next.\par

Shield released the exits one gallery at a time.\par

Even dispersal had doctrine.\par

Even leaving the Hall was a rite of being counted out safely from the place where one had been counted in.\par

The old Keeper’s chair remained empty on the dais for several minutes, dark wood, polished arms, a vacancy more eloquent than a body, and I watched it until two attendants came to take it away, not hurriedly, not reverently, but with the trained discretion of people removing a sign whose usefulness had expired.\par

Va-Raedin took my hand and bound the cut with a strip of pale cloth.\par

His fingers were efficient.\par

Not gentle.\par

Not ungentle.\par

Light had no use for either category when tightness would do.\par

“You will report to the eastern practice hall before first work,” he said.\par

“Do not eat heavily.\par

Do not arrive proud.”\par

“I am not proud,” I said.\par

Va-Raedin tied the cloth too tightly.\par

“You are fifteen.\par

Pride is your natural condition.”\par

I wanted to object, but objection required a path, and every path available led toward proving him right, so I kept silent, and he saw that too, of course he saw it, the little calculation passing across my face before discipline arrived to cover it, the boy not yet humble but already learning the uses of appearing so.\par

“Good,” he said.\par

“You have learned one thing already.”\par

There were men in the Mountain who could praise you and make you smaller.\par

Va-Raedin could correct you and leave you standing taller than you had any right to stand.\par

Across the Hall, Va-Sheva moved with the Shield wardens toward the northern passage, and my gaze followed before I could discipline it, past the dispersing crowd, past children being lifted, elders being steadied, clerks gathering tablets, wardens narrowing exits, until for one brief instant she looked back.\par

She smiled then.\par

Not as a woman smiles at a man.\par

I knew that, even in the foolish centre of myself that wanted to deny it.\par

She smiled as one who remembered a boy and approved that he had survived a difficult morning, and because kindness from Va-Sheva was never careless, never cheaply given, never false enough to comfort fully, it should have been enough.\par

It was not.\par

The appetite in me wanted rank disguised as tenderness.\par

It wanted the smile to cross a distance she had not crossed.\par

It wanted to steal more meaning than she had placed there.\par

Then she was gone into the Shield passage, and the arch swallowed her without ceremony, which was proper, because Shield did not decorate departures any more than it decorated threats.\par

I stood with my bandaged thumb at my side and my new name settling around me badly, like a robe cut for a body I had not yet grown, Sa-Kailan, Candidate of Light, spoken, cut, entered into service, and beneath the sting of Va-Raedin’s correction and the ache of Va-Sheva’s insufficient smile, pride rose in me, bright, undeniable, already looking for a lamp to justify itself.\par

I had been chosen.\par

That was what the boy knew.\par

I had been owned more precisely.\par

That is what Switzerland taught the man to say.\par

Faith had walked until the body nearly failed.\par

Balance had measured water before it became blessing.\par

Shield had watched the exits while others watched the altar.\par

Justice had written the event before memory could improve it.\par

Light, I thought then, must do something cleaner.\par

Even now I almost pity that thought.\par

The Hall emptied around me, lamps dimming one by one as attendants pinched the wicks, and above the dais the gravure darkened until the carved Mountain lost its edges and the Oasis disappeared into shadow, so that for a moment the old stone no longer showed refuge, promise, or history, only a blackened place where meaning had been, and I looked at that shadow until my eyes began to invent shapes inside it.\par

Then I turned toward the eastern passage and followed Va-Raedin into Light.\par

\ornament


\ornament
